%% ============================================================================
%% CAPÍTULO 5: EL MAP MATCHING Y EL ÍNDICE ESPACIAL
%% ============================================================================
\chapter{Contribución II: Introducción de un índice espacial para la ingestión eficiente de Dataframes}

\section{El Problema de Asignación Espacial a Gran Escala}
% GUÍA: Explicar el problema del "vecino más cercano" aplicado a crímenes (ej. un delincuente escapando entre esquinas) y por qué las búsquedas lineales o en grafos crudos fallan.

\section{Diseño del Índice Espacial mediante Geometría Computacional}
% GUÍA: Explicar el uso de la librería CGAL en C++ para particionamiento espacial. 
% Detallar la implementación basada en diagramas de Voronoi y Árboles K-D para búsquedas en tiempo logarítmico.

\section{Ingestión Eficiente de DataFrames}
% GUÍA: Cómo el algoritmo procesa miles de coordenadas GPS (DataFrames de crímenes) en bloque, minimizando el cuello de botella entre memoria y CPU.

\section{Análisis de Complejidad Computacional}
% GUÍA: Demostración teórica de la mejora en la notación Big-O (Ej. de O(N*M) a O(N log M)) gracias al índice espacial estructurado.